\chapter{El kernel \texorpdfstring{FOL$^=$}{FOL=}}
\label{cap:kernel}

Antes de poder decir «esta teoría no demuestra su propia consistencia» hay que tener una teoría, y
antes de tener una teoría hay que tener un lenguaje y una noción de demostración. Ese es el papel del
proyecto \modulo{FOL}, hermano de éste: lo que aquí se usa, y sólo eso, es lo que este capítulo
resume.

\begin{leccion}[Criterio de selección de este capítulo]
\modulo{FOL} es más grande que lo que este libro necesita: tiene semántica, solidez, completitud y
compacidad. La cadena de la incompletitud \textbf{no importa ninguno de esos módulos} — es
íntegramente sintáctica. Lo que sigue es el corte medido, no una visita guiada.
\end{leccion}

\section{El lenguaje}

\begin{matematica}[Términos y fórmulas]
Los términos son variables (índices de De Bruijn) o aplicaciones de un símbolo de función a una
lista de términos:
\[ t \;::=\; v_n \;\mid\; f(t_1,\dots,t_k). \]
Las fórmulas son
\[ A \;::=\; \bot \;\mid\; P(\bar t) \;\mid\; t \doteq u \;\mid\; A \Rightarrow B
   \;\mid\; \forall A \;\mid\; A \wedge B \;\mid\; A \vee B \;\mid\; \exists A. \]
La negación, la verdad y el bicondicional no son primitivas:
\[ \lnot A := A \Rightarrow \bot, \qquad \top := \lnot\bot, \qquad
   A \Leftrightarrow B := (A \Rightarrow B) \wedge (B \Rightarrow A). \]
\end{matematica}

\leanfrag{fol_Term}
\leanfrag{fol_Formula}

Nótese que los cuantificadores \textbf{no llevan variable}: \(\forall A\) liga la variable
\(\texttt{\#0}\) del cuerpo. Es la consecuencia visible de la decisión siguiente.

\section{De Bruijn, y el precio que se paga}

Con nombres de variable, la sustitución exige renombrar para evitar capturas, y toda prueba formal
arrastra hipótesis de frescura. Con índices no hay nombres que capturar: una variable es la distancia
al cuantificador que la liga. La captura desaparece — y aparece la aritmética de índices.

\begin{matematica}[Desplazamiento y sustitución]
El desplazamiento \(\uparrow^{c}\) incrementa las variables \emph{libres} por encima del corte
\(c\); la sustitución \(A[s/v]\) reemplaza \(v\) por \(s\) y decrementa lo que quedaba por encima:
\[
\uparrow^{c}(v_n) =
  \begin{cases} v_n & n < c\\ v_{n+1} & n \ge c\end{cases}
\qquad
v_n[s/v] =
  \begin{cases} s & n = v\\ v_{n-1} & n > v\\ v_n & n < v.\end{cases}
\]
Al cruzar un cuantificador ambas suben un nivel, y la sustitución debe desplazar lo que introduce:
\(\forall A [s/v] = \forall\bigl(A[\uparrow^{0}\!s\,/\,v{+}1]\bigr)\).
\end{matematica}

\leanfrag{fol_liftTerm}
\leanfrag{fol_substTerm}
\leanfrag{fol_substFormula}

\begin{leccion}[El precio, medido]
La última línea de \ident{substFormula} —el \ident{liftTerm 0 s} al entrar en un cuantificador— es
la fuente de casi toda la fricción técnica del proyecto. No es una impresión: contando las citas de
\modulo{FOL} en los 108 módulos de ROBINSON\texttt{++}, \textbf{tres lemas de conmutación
lift/subst acumulan 1\,316 usos}, y el más citado de todos, \ident{substTerm_liftTerm}, aparece
\textbf{978 veces}. Formalizar con De Bruijn no es más difícil que con nombres: es difícil en otro
sitio, y el sitio es éste.
\end{leccion}

\section{El cálculo \texorpdfstring{$\vdash$}{|-}}

\begin{matematica}[Deducción natural sobre listas de hipótesis]
\(\Gamma \vdash A\) se define inductivamente con las reglas usuales de deducción natural
—hipótesis, introducción y eliminación de \(\Rightarrow\), \(\wedge\), \(\vee\), \(\forall\),
\(\exists\)—, más \emph{ex falso}, debilitamiento, y las dos reglas de igualdad
\[
\frac{}{\Gamma \vdash t \doteq t}\quad(\text{refl})
\qquad
\frac{\Gamma \vdash t_1 \doteq t_2 \qquad \Gamma \vdash A[t_1]}{\Gamma \vdash A[t_2]}\quad(\text{subst}).
\]
La regla de \(\forall\)-introducción desplaza todo el contexto
(\(\Gamma\!\uparrow^{0} \vdash A\) da \(\Gamma \vdash \forall A\)): es la forma De Bruijn de exigir
que la variable no aparezca libre en las hipótesis.
\end{matematica}

Hay además una regla que no es estándar y conviene señalar: \ident{rewrite_at} permite aplicar una
regla local \emph{en una posición exacta} del árbol sintáctico, navegando con \ident{Pos} y
\ident{getAt?}. Es maquinaria de reescritura, no fuerza lógica añadida.

El teorema de deducción es aquí un teorema de una línea, porque el cálculo ya tiene
\ident{intro_impl} como constructor:

\leanfrag{fol_deduction_theorem}

\section{Las meta-reglas \texorpdfstring{$\omega$}{omega}: seis axiomas, y lo que cuestan}

Sobre \(\vdash\) hay una segunda capa, en \modulo{FOL/MetaRules.lean}, cuyas hipótesis son
\textbf{funciones de Lean} en vez de hipótesis en contexto. Seis de ellas son \ident{axiom}:
\ident{imp_intro}, \ident{gen}, \ident{raa}, \ident{dne}, \ident{or_elim} y \ident{ex_elim}.

\leanfrag{fol_imp_intro}
\leanfrag{fol_gen}

Que sean axiomas no es una comodidad: es que \textbf{no son derivables}, y el contraejemplo es
concreto. Tómese \(\Gamma = [\,]\), \(A = P\), \(B = Q\) átomos distintos. En Lean clásico
\(\lnot(\;[\,] \vdash P\;)\) es cierto, luego la hipótesis \(([\,]\vdash P) \to ([\,]\vdash Q)\) se
cumple \emph{vacuamente}; pero \([\,] \vdash P \Rightarrow Q\) no es derivable. Aceptar
\ident{imp_intro} como teorema de FOL sería, por tanto, aceptar algo falso. Lo mismo vale para
\ident{raa}, \ident{or_elim} y \ident{ex_elim}. La forma correcta a nivel objeto existe y se
prefiere en código nuevo: \ident{Derives.intro_impl}.

\ident{gen} es de otra naturaleza: es la \textbf{\(\omega\)-regla}. De «\(\Gamma \vdash A[n]\) para
todo término \(n\)» concluye \(\Gamma \vdash \forall A\). Es válida en el modelo estándar
\(\mathbb{N}\) y no es derivable en ningún sistema finitario.

\begin{muro}[La consecuencia que gobierna todo el libro]
Con la \(\omega\)-regla, \(\vdash\) tiene reglas de \textbf{premisas infinitas}: no es recursivamente
enumerable. Y por Tarski, ninguna fórmula \(\Prov\) puede satisfacer
\(\vdash \Prov(\codigo{\varphi}) \Leftrightarrow \;\vdash \varphi\) para ese \(\vdash\).
Por eso el proyecto construye, \emph{en paralelo}, un cálculo de Hilbert \textbf{finitario}
\(\Prf\), y toda la aritmetización se hace sobre él. Los dos cálculos y su puente son el capítulo siguiente;
aquí basta con saber que hay dos, y por qué tiene que haberlos.
\end{muro}

\section{Cuatro relaciones de derivabilidad, no dos}

De la observación anterior no sale «un cálculo más»: salen \textbf{cuatro} relaciones distintas, y
lo que las separa —si arrastran contexto y si son recursivamente enumerables— es exactamente lo que
no se puede confundir. Sus tipos en Lean lo dicen:

\leanfrag{fol_Derives}
\leanfrag{Prf_tipo}
\leanfrag{PrfH_tipo}

\begin{matematica}[Las cuatro, y sus puentes]
\[
\begin{array}{lccl}
\Gamma \vdash A & \text{con contexto} & \text{\textbf{no} r.e.} & \text{cálculo }\omega\\
\Prf_0\, A      & \text{sin contexto} & \text{r.e.} & \text{Hilbert intuicionista}\\
\Prf\, A        & \text{sin contexto} & \text{r.e.} & \text{Hilbert clásico: \textbf{el que se aritmetiza}}\\
\mathsf{PrfH}\ \Gamma\ A & \text{con contexto} & \text{r.e.} & \text{variante contextual}
\end{array}
\]
\[
  \Prf_0\varphi \to \Prf\varphi \to \bigl(\forall\Gamma,\ \mathsf{PrfH}\,\Gamma\,\varphi\bigr),
  \qquad \Prf\varphi \to \text{\ident{axioms}} \vdash \varphi .
\]
\end{matematica}

\leanfrag{prf0_to_derives}
\leanfrag{prf_to_derives}
\leanfrag{prf_deduction}

Los dos puentes hacia $\vdash$ se diferencian en \textbf{un solo símbolo}, y está medido: el primero
no invoca \emph{ninguna} meta-regla $\omega$; el segundo usa \ident{dne} exactamente una vez. Es la
frontera intuicionista/clásica, aislada en un punto.

\begin{muro}[El recíproco no existe, y tiene que no existir]
$\text{\ident{axioms}} \vdash \varphi \to \Prf\,\varphi$ es \textbf{falso}. Si valiera, $\vdash$
sería recursivamente enumerable y Tarski cerraría el paso: ningún $\Prov$ podría representar su
demostrabilidad. La asimetría no es un defecto de la formalización — \emph{es} la razón de que haya
dos cálculos.
\end{muro}

\section{El toolkit De Bruijn que de verdad se usa}

De los diecinueve lemas de \modulo{FOL/Theorems/Eq.lean}, el proyecto vive esencialmente de cinco.
Ordenados por número de citas medidas en ROBINSON\texttt{++}:

\leanfrag{fol_substTerm_liftTerm}
\leanfrag{fol_derive_eq_trans}
\leanfrag{fol_substTerm_liftLift}
\leanfrag{fol_derive_eq_symm}
\leanfrag{fol_liftTerm_comm_zero}

\begin{center}\small
\begin{tabular}{lr}
\toprule
lema & citas en \modulo{ROBINSON\_PlusPlus/} \\
\midrule
\ident{substTerm_liftTerm} & 978 \\
\ident{derive_eq_trans}    & 429 \\
\ident{substTerm_liftLift} & 312 \\
\ident{derive_eq_symm}     & 57 \\
\ident{liftTerm_comm_zero} & 26 \\
\bottomrule
\end{tabular}
\end{center}

\noindent Los dos primeros dicen, respectivamente, que \emph{sustituir deshace desplazar} y que la
igualdad derivada es transitiva. Que un proyecto de metamatemática se apoye casi por completo en
esas dos cosas es, en sí, el resumen honesto de en qué consiste formalizar.

\section{Un axioma que era falso}

\modulo{FOL/Theorems/Quantifiers.lean} declaraba en su día:
\[ \texttt{substFormula } v\ t\ (\texttt{liftFormula } (v{+}1)\ f) \;=\; f
   \qquad\text{para \emph{cualquier} } t. \]

\begin{refutado}[El contraejemplo, en tres líneas]
Sea \(f = P(\texttt{\#0})\), \(v = 0\), \(t = \texttt{zero}\). Entonces
\(\uparrow^{1}\! f = P(\texttt{\#0})\) —la variable \(0 < 1\) no se desplaza— y
\(P(\texttt{\#0})[\texttt{zero}/0] = P(\texttt{zero}) \neq f\).
\end{refutado}

El enunciado correcto exige que el término sustituido sea \emph{exactamente} la variable, y entonces
es un teorema, demostrable por inducción estructural:

\leanfrag{fol_subst_lift_cancel}

\begin{leccion}[Por qué esto abre el libro y no lo cierra]
Es el mismo modo de fallo que la Parte~IV cuenta en grande: un enunciado plausible, aceptado como
axioma, que resulta ser falso —y que mientras estuvo ahí hacía \emph{unsound} todo lo que lo citaba.
Aquí costó un cambio de enunciado; en el capítulo~\ref{cap:numerales} costó 31 módulos. La
diferencia no es de naturaleza, sino de cuánto se tardó en mirar.
\end{leccion}

\section{Lo que hay en FOL y este libro no usa}

Para que el alcance quede declarado y no haya que deducirlo:

\begin{center}\small
\begin{tabularx}{\linewidth}{@{}>{\raggedright\arraybackslash}p{0.30\linewidth}
                              >{\raggedright\arraybackslash}X
                              >{\raggedright\arraybackslash}p{0.14\linewidth}@{}}
\toprule
módulo de \modulo{FOL} & contenido & ¿se usa? \\
\midrule
\modulo{FOL/FOL.lean} & sintaxis, De Bruijn, \ident{Derives} & sí, todo \\
\modulo{FOL/MetaRules.lean} & las 6 meta-reglas \(\omega\) & sí \\
\modulo{FOL/Theorems/Eq.lean} & toolkit De Bruijn e igualdad & sí, cinco lemas \\
\modulo{FOL/Deduction.lean} & teorema de deducción & sí \\
\modulo{FOL/Theorems/}: \ident{Impl}, \ident{Neg}, \ident{Derived}, \ident{Quantifiers}
  & lógica proposicional y de cuantificadores & parcialmente \\
\modulo{FOL/}: \ident{Semantics}, \ident{Soundness}, \ident{Completeness}, \ident{Compacity}
  & modelos, solidez, completitud, compacidad & \textbf{no} \\
\bottomrule
\end{tabularx}
\end{center}

\begin{observacion}
La mitad semántica de \modulo{FOL} no la importa \textbf{ningún} módulo de
ROBINSON\texttt{++} (verificado: cero \ident{import}). Esto importa porque
\modulo{FOL/Completeness.lean} declara cinco \ident{axiom} propios —entre ellos una enumeración de
fórmulas y el lema de extensión de Henkin— que, por tanto, no pueden entrar en la base de ningún
teorema de este libro. Y por la misma razón conviene decirlo explícitamente: \modulo{FOL} declara
\textbf{trece} \ident{axiom} en total, de los que la base sancionada de este libro admite
\textbf{seis}. Los otros siete están fuera de la cadena, y \printaxioms\ es lo que lo comprueba.
\end{observacion}
